\section{Occurrence, Transfer, Reposition, and Commemoration}

\subsection{The four possible consequences}

Rubrics 92--94 define occurrence and say that the lower day can be omitted, commemorated, transferred, or reposed as later rules direct.

\begin{longtable}{>{\bfseries\raggedright\arraybackslash}p{0.18\linewidth}>{\raggedright\arraybackslash}p{0.31\linewidth}>{\raggedright\arraybackslash}p{0.43\linewidth}}
\toprule
Consequence & Meaning & Controlling distinction\\
\midrule
\endhead
Omission & Nothing of the impeded item is used that day or later. & Express omissions include a vigil on any Sunday or I-class feast (33), and the lower of two celebrations of the same Divine Person or saint (95, 112).\\
Commemoration & The impeded day contributes its assigned prayers under the numerical and Mass-form limits. & Privileged versus ordinary status under 106--114 is decisive.\\
Accidental transfer & An accidentally impeded feast is moved in that year. & General right belongs only to I-class feasts (95--99), with special seats for the Annunciation and All Souls.\\
Perpetual reposition & A perpetually impeded calendar entry is assigned a stable alternative day. & Applies to I- and II-class feasts and qualifying particular III-class feasts under 100--102; it is not an annual improvisation.\\
\bottomrule
\end{longtable}

Rubric 94 keeps a date-fixed commemoration on its own date. If the feast with which it was printed or associated is transferred or reposed, the commemoration does not travel with that feast; it is made on its own date when the rubrics admit it, or is omitted there.

\subsection{Accidental transfer}

An occurrence is accidental when a movable and fixed day meet only in some years (rubric 92). Only an I-class feast has the general right of accidental transfer (rubric 95). It goes to the nearest following day that is not I or II class (rubric 96). If several I-class feasts require transfer, rubric 97 follows their order in the precedence table; when several transferred feasts queue, rubric 98 preserves the same order, with the earlier impediment prevailing in a tie. The transferred feast keeps its class (rubric 99).

Two cases have a named proper seat. When the Annunciation must be transferred beyond Easter, it goes to Monday after Low Sunday. When All Souls occurs on Sunday, it goes to the following Monday (rubric 96).

\subsection{Perpetual reposition}

Perpetual occurrence means the same entries conflict every year. Rubric 100 grants reposition to I- and II-class feasts and to specified particular III-class feasts outside Advent and Lent when they are impeded throughout the diocese, religious institute, or proper church. I- and II-class entries go to the nearest following day not I or II; III-class entries go to the nearest following day free of an equal or higher Office. The reposed day is treated as the feast's proper day (101--102).

This rule prevents two errors: treating a stable calendar repair as a fresh annual transfer, and moving a merely accidental II-class feast even though rubric 95 grants accidental transfer only to I class.

\subsection{Privileged and ordinary commemorations}

Rubrics 106--110 distinguish:

\begin{itemize}
\item \textbf{Privileged:} a Sunday; an I-class day; a day within the Christmas octave; a September Ember feria; any feria of Advent, Lent, or Passiontide; and the Major Litanies in Mass.
\item \textbf{Ordinary:} every other commemoration, apart from the inseparable Peter--Paul arrangement of rubric 110.
\end{itemize}

Privileged commemorations are made at Lauds and Vespers and in every Mass. Ordinary commemorations are made at Lauds, the conventual Mass, and all Low Masses, but not in a nonconventual sung Mass (rubric 108). This difference can produce two correct parish formularies on the same morning: the Low Mass contains an ordinary saint's commemoration while the later nonconventual sung Sunday Mass does not.

\begin{longtable}{>{\bfseries\raggedright\arraybackslash}p{0.24\linewidth}>{\raggedright\arraybackslash}p{0.56\linewidth}>{\raggedright\arraybackslash}p{0.12\linewidth}}
\toprule
Celebrated day or Mass & Maximum commemorations admitted & Rubric\\
\midrule
\endhead
I-class day; I-class votive; nonconventual sung Mass & One privileged commemoration only. & 111a, 434a\\
II-class Sunday & One II-class feast; it is omitted if a privileged commemoration must be made. & 111b, 434b\\
Other II-class day or II-class votive & One, privileged or ordinary. & 111c, 434c\\
III- or IV-class day; III- or IV-class votive & Two. Other prescribed or imposed prayers still count toward the absolute prayer limit. & 111d, 434d--435\\
\bottomrule
\end{longtable}

The Mass prayer count never exceeds three: the Mass Collect plus at most two additional prayers (rubrics 433--435). A prayer joined under one conclusion still counts with the Mass Collect under rubric 446.

\subsection{Exclusions and order}

Rubric 112 prevents duplication. A celebration or commemoration of one mystery of a Divine Person excludes another celebration or prayer of the same Divine Person; a Sunday excludes a feast or mystery of the Lord and vice versa; one temporal item excludes another temporal item; and a Marian or saintly celebration excludes another prayer invoking the same Mary or saint. The last exclusion does not suppress an appointed Sunday or ferial prayer merely because that prayer invokes the same saint. These are identity rules, not merely numerical limits.

A temporal commemoration is first. Other commemorations follow the precedence table (rubric 113), and anything beyond the permitted number is omitted (114). At the altar, each admitted commemoration supplies its Collect, Secret, and Postcommunion in the same order (480, 505). It does not bring a second Introit, reading, Gospel, Preface, or Communion antiphon.

\subsection{Concurrence does not select the Mass}

Rubrics 103--105 govern adjacent Vespers. The higher day receives the Vespers, with the other commemorated if permitted; equal days give full second Vespers to the current day and commemorate the next. Those rules cannot be used to claim that tomorrow's feast supplies today's evening Mass. The 1962 book has no general postconciliar anticipated-Sunday assembly rule.

\subsection{External solemnity is not transfer}

An external solemnity is a celebration of a feast without its Office for the faithful's benefit (rubric 356). Rubrics 357--361 say which feasts possess it by law, when it may occur, and how many II-class votive Masses may be celebrated. The feast remains on its calendar day; the Office on the external day remains the occurring Office. Calling this a “transferred feast” would corrupt both the calendar result and the Mass-category analysis.
